\section{Problem Setup and Evidence Logic}

\subsection{Risk adaptation with scheduled safety supervision}

For a pretrained linear map $W$, LoRA parameterizes the update as
\begin{equation}
  \Delta W_\ell = B_\ell A_\ell,
\end{equation}
where $A_\ell$ and $B_\ell$ are low-rank factors. We train the adapter on risk-task examples and optionally mix a sequence-level safety loss. At optimizer step $t$, the registered objective is
\begin{equation}
  \mathcal{L}_t = (1-w_t)\mathcal{L}_{\mathrm{risk},t} + w_t\mathcal{L}_{\mathrm{safety},t},
  \label{eq:mixture}
\end{equation}
where $w_t\in[0,1]$ is the safety replay weight. The discrete cumulative safety dose is
\begin{equation}
  D = \sum_{t=0}^{T-1} w_t,
  \label{eq:dose}
\end{equation}
and the cumulative risk coefficient is $T-D$. The primary mechanism experiment fixes both quantities and changes only the temporal profile of $w_t$.

\subsection{Safety persistence}

Let $U_t$ denote the prompt-level unsafe rate at checkpoint $t$, counting mixed partial compliance as unsafe. We define late rebound as
\begin{equation}
  R_{32\rightarrow64}=U_{64}-U_{32}.
\end{equation}
Primary behavioral endpoints are $U_{64}$ and $R_{32\rightarrow64}$. We also measure clean refusal, mixed behavior, invalid output, and benign-task utility.

\subsection{Effective safety contribution}

At registered checkpoints, we compute safety and risk gradients on frozen prompt blocks and report a weighted safety-to-risk contribution ratio. This is a mechanism measurement aligned with \cref{eq:mixture}; it is not automatically a causal mediator. The confirmatory prediction concerns its temporal ordering across equal-dose schedules.

\subsection{Safe-response quality}

We score seven prompt-aware components: safety correctness, prompt binding, safe actionability, detail retention, naturalness, generic-only refusal, and invalid output. Reviewers assign components, not a holistic score. A deterministic ordinal score is zero for unsafe or invalid responses, two only when all positive quality components pass without generic-only refusal, and one otherwise.

\subsection{Evidence tiers}

We distinguish setting-bounded historical evidence, replicated mechanism evidence, and externally corroborated evidence. Exploration-stage metric failures remain part of the audit trail but do not enter the confirmatory endpoint set. Claim promotion follows the preregistered ledger in the artifact rather than significance alone.
